\begin{table}[H]
\centering
\caption{Standardized Effect Sizes for Main Outcomes}
\label{tab:sde}
\begin{threeparttable}
\footnotesize
\begin{tabular}{llccccc}
\toprule
Outcome & $\hat{\beta}$ & SE & SD($Y$) & SDE & SE(SDE) & Classification \\
\midrule
\multicolumn{7}{l}{\textit{Panel A: Pooled}} \\
Employment (log) & -0.0612 & 0.0552 & 1.1546 & -0.0530 & 0.0478 & Moderate negative \\
New hires (log) & -0.0120 & 0.0594 & 1.1175 & -0.0107 & 0.0532 & Small negative \\
Earnings (log) & -0.0078 & 0.0219 & 0.2970 & -0.0263 & 0.0739 & Small negative \\
\midrule
\multicolumn{7}{l}{\textit{Panel B: Heterogeneous (by worker ethnicity)}} \\
Employment, Hispanic & -0.0719 & 0.0702 & 1.2974 & -0.0554 & 0.0541 & Moderate negative \\
Employment, Non-Hispanic & 0.0151 & 0.0461 & 1.2925 & 0.0117 & 0.0357 & Small positive \\
\bottomrule
\end{tabular}
\begin{tablenotes}[flushleft]
\scriptsize
\item \textit{Notes:} \textbf{Country:} United States. \textbf{Research question:} Do state-level constitutional Right-to-Farm amendments, which immunize concentrated animal feeding operations from nuisance litigation and local zoning, increase county-level animal production employment? \textbf{Policy mechanism:} Constitutional RTF amendments create a legal shield that eliminates or severely restricts the ability of neighbors, local governments, and environmental groups to sue large-scale animal operations for nuisance, odor, water contamination, or zoning violations---effectively removing a major regulatory cost and litigation risk for facility operators. \textbf{Outcome definition:} Quarterly county-level beginning-of-quarter employment count in NAICS 112 (Animal Production and Aquaculture) from the Census Bureau's Quarterly Workforce Indicators, measured in logs. \textbf{Treatment:} Binary indicator for state adoption of a constitutional RTF amendment or major legislative RTF strengthening, with staggered timing across six states between 2012 and 2021. \textbf{Data:} QWI county-quarter panel, 2005Q1--2024Q4, covering 14 states (6 treated, 8 agricultural controls). \textbf{Method:} Staggered difference-in-differences with TWFE estimator, county and quarter fixed effects, standard errors clustered at the state level (14 clusters), supplemented by wild cluster bootstrap for finite-sample inference. \textbf{Sample:} County-quarters with non-suppressed NAICS 112 employment in 6 RTF-strengthening states (ND, NC, MO, IA, GA, TX) and 8 agricultural control states (KS, NE, WI, TN, VA, SD, KY, MT) without RTF changes during the sample period. SDE $= \hat{\beta} / \text{SD}(Y)$ where SD($Y$) is the pre-treatment standard deviation of log employment. Classification refers to magnitude, not statistical significance: Large ($|$SDE$| > 0.15$), Moderate ($0.05$--$0.15$), Small ($0.005$--$0.05$), Null ($< 0.005$).
\end{tablenotes}
\end{threeparttable}
\end{table}
